% 当前英文版第三至第五章的中文审阅稿。
% 基准源文件：Chapter3/ch_model.tex、Chapter4/ch_framework.tex、Chapter5/ch_eval.tex
% 截止日期：2026-11-05。保留章节结构、关键公式、主要表格和全部结论；
% 为方便审阅，对英文原文中重复的机制说明作了合并，不作为最终提交版本。
\documentclass[11pt,UTF8,fontset=none]{ctexrep}
\setCJKmainfont{Noto Serif CJK SC}[BoldFont={Noto Serif CJK SC SemiBold}]
\setCJKsansfont{Noto Sans CJK SC}[BoldFont={Noto Sans CJK SC Medium}]
\usepackage[a4paper,top=2.5cm,bottom=2.5cm,left=2.8cm,right=2.8cm,headheight=16pt]{geometry}
\usepackage{amsmath,amssymb,bm}
\usepackage{booktabs,tabularx,array}
\usepackage{graphicx,float}
\usepackage{siunitx}
\usepackage{xcolor}
\usepackage{setspace}
\usepackage{fancyhdr}
\usepackage[hidelinks]{hyperref}
\sisetup{detect-all}

\providecommand{\dJ}{\Delta J}
\providecommand{\mB}{m_B}
\providecommand{\mP}{m_P}
\providecommand{\mT}{m_T}
\providecommand{\Jb}{J_B}
\providecommand{\Jt}{J_T}
\providecommand{\svec}{\bm{s}}
\providecommand{\cvec}{\bm{c}_{xy}}
\providecommand{\rp}{\bm{r}_p}
\providecommand{\rc}{\bm{c}}
\providecommand{\Tphys}{T_{\mathrm{phys}}}
\providecommand{\tauphys}{\bm{\tau}_{\mathrm{phys}}}
\providecommand{\ev}{\bm e_3}
\DeclareMathOperator*{\argmin}{arg\,min}

\pagestyle{fancy}
\fancyhf{}
\fancyhead[L]{第三至第五章中文审阅稿}
\fancyhead[R]{\leftmark}
\fancyfoot[C]{\thepage}
\setlength{\parindent}{2em}
\setlength{\parskip}{0.35em}
\setlength{\emergencystretch}{3em}
\hypersetup{pdftitle={学位论文第三至第五章中文审阅稿}}
\graphicspath{{src/paper/thesis/}}

\title{学位论文第三至第五章中文审阅稿\\[4pt]
\large 突变载荷下基于移动时域估计的四旋翼非线性模型预测控制}
\author{}
\date{基于 2026 年 11 月 5 日英文工作稿}

\begin{document}
\maketitle
\tableofcontents

\chapter*{阅读说明}
\addcontentsline{toc}{chapter}{阅读说明}
本稿用于快速检查当前英文论文第三至第五章的技术内容与论证逻辑。它保留当前版本的章节结构、关键公式、主要实验表格、样本量、统计结论和局限；为提高中文可读性，将英文原文中多次出现的同一机制说明合并。公式符号与英文版一致。

当前英文工作稿仍包含一项尚未完成的内容：全状态估计器 F 与速度--质量降阶估计器 R 的 96 架闭环配对实验仍在运行。因此，第五章相应位置只翻译已经写入英文源文件的回放、四架冒烟实验和预注册设计，不虚构闭环结果。

% =====================================================================
\setcounter{chapter}{2}
\chapter{建模与参数结构}

四旋翼抓取包裹后不仅质量增加，滚转和俯仰惯量也会显著增大。在本文平台上，悬挂于旋翼平面下方约四分之一米的 \SI{0.30}{\kg} 包裹只使总质量增加约 \SI{15}{\percent}，却使滚转和俯仰惯量增加到空机的两倍以上。因此，只补偿质量无法解释内环带宽下降。

本章的核心结构结论是：当载荷位于推力轴上时，质量变化和惯量变化可由一个在线标量统一描述。估计器通过平动通道辨识复合质量 $\mT$，随后利用刚体关系和夹爪的名义竖直力臂计算滚转/俯仰惯量增量。估计器和控制器都不接收任务给定的载荷质量，估计器也不接收吸附或释放通知。

\section{系统定义与建模假设}

世界坐标系 $\mathcal W$ 采用东--北--天（ENU），重力沿 $-z$；机体系 $\mathcal B$ 采用前--左--上（FLU），原点位于旋翼盘几何中心，而不是复合质心。推力轴通过该原点。状态与输入为
\begin{equation}
x=[\bm p;\bm v;\bm q;\bm\omega]\in\mathbb R^{13},\qquad
u=[T;\bm\tau]\in\mathbb R^4,
\end{equation}
其中 $\bm p,\bm v$ 在世界系表示，$\bm q$ 为机体系到世界系的单位四元数，$\bm\omega$ 为机体系角速度，$T$ 为总推力，$\bm\tau$ 为机体系力矩。记空机、载荷及复合质量分别为 $\mB,\mP,\mT=\mB+\mP$。

模型采用三项假设：空机质心与机体原点重合；载荷视为刚性连接在推力轴上的质点，$\rp=[0,0,r_z]^\top$ 且 $r_z<0$；旋翼推力满足经过标定的静态平方律。抓取前通过相对位置对准载荷，剩余水平偏移作为模型近似处理。

\subsection{姿态表示}

采用 Hamilton 约定、标量在前的单位四元数
\begin{equation}
\bm q=[q_w,q_x,q_y,q_z]^\top,\qquad \|\bm q\|=1.
\end{equation}
对应旋转矩阵 $R(\bm q)$ 将机体系向量旋转到世界系。总推力的竖直分量比例为
\begin{equation}
\cos\phi\cos\theta=\ev^\top R(\bm q)\ev=1-2(q_x^2+q_y^2),
\end{equation}
实现中直接由四元数计算右端。

\section{带刚性载荷的四旋翼动力学}

运动学为
\begin{equation}
\dot{\bm p}=\bm v,\qquad
\dot{\bm q}=\tfrac12\,\Xi(\bm q)\bm\omega,
\end{equation}
其中 $\bm q^\top\Xi(\bm q)=0$，因而连续时间动力学保持四元数模长。运动学不包含质量或惯量，载荷只通过动力学进入模型。

平动动力学为
\begin{equation}
\dot{\bm v}=\frac{1}{\mT}\left(R(\bm q)\ev T-k_d\bm v\right)-g\ev.
\end{equation}
质量只以推力到加速度的标量增益 $1/\mT$ 出现。若推力由独立于控制模型的测量获得，普通近悬停飞行已经能够提供质量信息。投影到竖直方向并忽略竖直阻力得到
\begin{equation}
\mT(g+a_z)=T\cos\phi\cos\theta,
\end{equation}
该关系用于构造估计器初值。

转动动力学为
\begin{equation}
\dot{\bm\omega}=\Jt^{-1}\left(\bm\tau-\bm\omega\times\Jt\bm\omega\right).
\end{equation}
惯量和力矩必须关于同一点定义；自由飞行刚体应取复合质心。对于轴上载荷，复合质心仍位于推力轴上，所以总推力不产生关于复合质心的力矩。

\section{载荷引起的质量--惯量耦合}

定义复合质量和约化质量
\begin{equation}
\mT=\mB+\mP,\qquad
\mu=\frac{\mB\mP}{\mT}.
\end{equation}
吸附后复合质心移动到
\begin{equation}
\rc=\frac{\mP}{\mT}\rp.
\end{equation}
平行轴定理必须关于这个新质心同时计算空机和载荷两部分。对滚转或俯仰轴，合并两项可得
\begin{align}
\dJ&=\mB\left(\frac{\mP}{\mT}|r_z|\right)^2
    +\mP\left(\frac{\mB}{\mT}|r_z|\right)^2\nonumber\\
&=\frac{\mB\mP}{\mT}r_z^2=\mu r_z^2.
\end{align}
因此完整惯量为
\begin{equation}
\Jt=\Jb+\operatorname{diag}(\mu r_z^2,\mu r_z^2,0).
\end{equation}
惯量增量在滚转和俯仰方向相同、为对角形式，轴上质点对偏航惯量的增量为零。若错误地使用 $\mP r_z^2$，在本文重载条件下会高估约 \SI{14.5}{\percent}。

惯量比
\begin{equation}
\kappa=\frac{J_{xx}+\dJ}{J_{xx}}
\end{equation}
描述相同力矩下角加速度下降的倍数。取 $\mB=\SI{2.0643}{\kg}$、$J_{xx}=\SI{0.0142}{\kg\meter\squared}$、$r_z=\SI{-0.27}{\meter}$，得到
\begin{center}
\begin{tabular}{lccc}
\toprule
载荷 & $\mu$ [kg] & $\dJ$ [\si{\kg\meter\squared}] & $\kappa$\\
\midrule
\SI{0.15}{\kg} & 0.140 & 0.0102 & 1.72\\
\SI{0.30}{\kg} & 0.262 & 0.0191 & 2.34\\
\bottomrule
\end{tabular}
\end{center}
重载只使质量增加约 \SI{15}{\percent}，却使滚转/俯仰惯量超过空机的两倍。任务难度因此同时由载荷质量和竖直力臂决定。

\section{单参数估计结构}

载荷质量会随抓取和释放变化，并通过平动方程中的 $1/\mT$ 可观；竖直力臂则是夹爪几何属性，在配送飞行中经转动通道弱激励。因此在线只估计复合质量，并定义
\begin{equation}
\mP=(\mT-\mB)^+,
\qquad
\boxed{\dJ(\mT)=\frac{\mB(\mT-\mB)^+}{\mT}r_z^2}.
\end{equation}
当质量估计回到空机值时，惯量增量连续归零，不需要吸附或释放事件去切换。

\subsection{质量由哪个通道辨识}

理论上，显著的惯量变化会让滚转/俯仰运动包含质量信息；但本文平台的转动模型存在三类误差：标定惯量与仿真模型描述差异较大，旋翼平面竖直偏置引起的速度相关阻力矩未建模，空机重构力矩还带有小的常值偏置。若允许优化器通过质量直接改变转动模型，它会用质量吸收这些误差，开发实验中曾造成大幅质量高估。

部署的估计器因此完全不积分转动动力学。它只保留世界系速度与复合质量，把实测姿态作为已知输入。质量只进入可信的“推力--质量--重力”平动关系，估计器内部不再出现惯量、力矩或力臂。

第五章前半部分的实验使用较早的14状态全状态估计器。该估计器积分姿态和角速度，但在一个估计窗口内把惯量增量作为由上次质量解计算的已知常数：
\begin{equation}
\dJ^{\mathrm{lag}}=
\frac{\mB(\hat m^- -\mB)^+}{\hat m^-}r_z^2,qquad
\left.\frac{\partial\dot{\bm\omega}}{\partial\mT}\right|_{\mathrm{MHE}}=0.
\end{equation}
这种滞后处理去除了质量对转动动力学的直接灵敏度，但角速度残差仍可通过估计姿态影响推力方向，上一窗口的质量也会反馈到下一窗口。降阶估计器同时消除了这两条间接路径。

\subsection{估计量、派生量与固定量}

唯一的在线决策变量是复合质量。载荷质量和滚转/俯仰惯量由它派生；$r_z$ 是夹爪名义几何；水平质心偏移固定为零；空机质量、惯量和旋翼系数通过标定获得。本文的“无载荷先验”主张准确含义是：估计器和控制器不接收某次任务的载荷质量，估计器不接收吸附/释放通知。系统仍使用两个平台级常数：夹爪名义力臂和 \SI{0.30}{\kg} 载荷包线，后者用于吸附阶段的惯量保护下限。

不单独估计惯量，是因为它在近悬停配送飞行中弱可观，而且只能由误差最大的转动通道辨识；不估计水平偏移，是因为当前工况通过对准保证居中，偏移释放后又失去定义；使用名义 $r_z$ 是因为它属于夹爪设计而不是某个包裹的信息，但必须承认其误差以平方形式传入 $\dJ$。

\section{实测输入重构与模型接口}

估计器使用的推力和力矩由实测旋翼转速重构，而不是读取控制器指令。对本文 X 型布局，
\begin{equation}
\begin{bmatrix}\Tphys\\\tauphys\end{bmatrix}
=M\begin{bmatrix}F_1\\F_2\\F_3\\F_4\end{bmatrix},
\qquad F_i=k_f\omega_i^2,
\end{equation}
其中 $M$ 的各行分别计算总推力、滚转力矩、俯仰力矩和偏航反扭矩。转速以 \SI{250}{\hertz} 采样，在每个 \SI{10}{\hertz} 估计区间内平均。部署的降阶估计器只使用总推力，全状态估计器使用完整重构力/力矩。

不能使用指令推力：若控制器按假定质量 $\tilde m$ 给出悬停推力 $T_{\rm cmd}=\tilde m g$，将它代回模型必然得到 $m=T_{\rm cmd}/g=\tilde m$，估计器只会复述控制器原有假设。实测旋翼转速位于分配器和电机动力学之后，不含估计状态，从而打破这条循环。

系统采用外环 NMPC 与 PX4 内环串联。预测力矩只存在于 NMPC 求解器内部；PX4 实际接收总推力和由预测轨迹插值得到的机体角速度设定值；估计器看到的是根据实际转速重构的物理力/力矩。这样既保留 PX4 的高带宽速率环和混控器，也使估计器使用“执行了什么”而不是“计划执行什么”。

\section{建模假设与局限}

轴上质点不产生推力偏心力矩、约化质量公式以及零偏航惯量增量在理想质点假设下精确成立。其余关系受以下近似限制：
\begin{itemize}
\item 水平偏移设为零。\SI{0.03}{\meter} 对准误差与 \SI{0.30}{\kg} 载荷产生约 \SI{0.09}{\newton\meter} 配平力矩，主要由 PX4 速率环积分项承担；相应惯量误差约为增量的 \SI{0.4}{\percent}。
\item 竖直力臂取固定的 \SI{-0.27}{\meter}，实际约为 \SIrange{0.25}{0.30}{\meter}。因此重载惯量增量相对模型约有 $-14$ 至 $+23$\% 的逐架误差。
\item 将包裹视为质点，忽略其自身中心惯量，量级约为滚转/俯仰增量的 \SI{2}{\percent}。
\item 空机转动模型近似，存在惯量、旋翼阻力矩和常值力矩偏置误差。这是质量与转动通道隔离的主要原因。
\item 全状态估计器的惯量比质量滞后一个 \SI{0.1}{\second} 求解周期，并忽略转移期间的 $\dot J\omega$ 项。
\item 旋翼模型为静态平方律，快速平移中的诱导流会破坏该模型。$k_f$ 的相对误差会一比一传入质量估计，平动数据无法同时区分推力比例和质量。
\end{itemize}

\section{本章小结}

轴上刚性载荷的滚转/俯仰惯量增量等于约化质量乘以竖直力臂平方。因此，只需在线辨识复合质量，惯量可以代数派生并随质量自动归零。质量通过由实测旋翼转速重构的推力和平动速度响应辨识。部署的降阶估计器不积分转动动力学，实测姿态只作为已知输入；实验阶段的全状态估计器则用滞后一拍的惯量参数实现较弱的隔离。下一章给出估计器、控制器和释放生命周期的具体数值形式。

% =====================================================================
\chapter{MHE--NMPC 框架}

第三章说明“估计什么”，本章说明“如何实现”。系统包含部署的速度--质量降阶估计器 R、用于既有实验批次的14状态全状态估计器 F、完整的 NMPC 预测模型、连续载荷接口、内环调度以及抓取--释放生命周期。估计器不接收吸附或释放通知；控制模型连续跟随估计；只有在持续证据成立后，控制器才接受一次释放并完成任务状态转换。

\section{闭环架构}

\begin{figure}[H]
\centering
\includegraphics[width=0.96\textwidth]{Figures/fig1_architecture_massonly_zh.pdf}
\caption{系统架构与信息流。估计链只消费里程计与旋翼转速；健康门控后的原子载荷帧同时驱动 NMPC 模型和机体角速度指令调度。}
\end{figure}

系统有三层：Gazebo/PX4/可分离夹爪构成物理与飞控层，PX4 速率控制器构成内环，ROS 中的 MHE 和 NMPC 构成外环。估计器只读里程计和旋翼转速，不读控制器计划输入、夹爪事件或仿真关节真值；仿真载荷真值只用于离线评分。

估计器发布一个原子帧，包含质量、由质量派生的惯量、无载荷证据和健康状态。NMPC 与内环调度读取同一帧，避免质量和惯量通过不同话题传播而形成不一致模型。MHE 以 \SI{10}{\hertz}、\SI{2}{\second} 窗口运行；NMPC 以 \SI{20}{\hertz}、\SI{1}{\second} 预测时域运行；设定值以 \SI{50}{\hertz} 发布。

\section{估计器与控制器模型}

\begin{table}[H]
\centering\small
\caption{同一对象的三种模型。}
\begin{tabularx}{\textwidth}{p{0.15\textwidth}XXX}
\toprule
 & \textbf{降阶估计器 R} & \textbf{全状态估计器 F} & \textbf{NMPC 预测模型}\\
\midrule
状态 & $[\bm v;\mT]\in\mathbb R^4$ & $[\bm p;\bm v;\bm q;\bm\omega;\mT]\in\mathbb R^{14}$ & $[\bm p;\bm v;\bm q;\bm\omega]\in\mathbb R^{13}$\\
已知输入 & 实测姿态、$\Tphys$ & $\Tphys,\tauphys$、滞后 $\dJ$ & 参考、$\mT,\dJ,\hat{\bm d}$\\
动力学 & 仅平动 & 完整平动与转动 & 完整预测模型\\
残差 & 三维速度 & 十三维状态 & 跟踪与输入偏差\\
频率/时域 & \SI{10}{\hertz}, $N=20$ & 相同 & \SI{20}{\hertz}, $N=20$\\
\bottomrule
\end{tabularx}
\end{table}

\subsection{部署的速度--质量估计器}

降阶状态和已知参数为
\begin{equation}
x_R=[\bm v;\mT]\in\mathbb R^4,qquad
p_{R,j}=[\tilde{\bm q}_j;T_{{\rm phys},j}]\in\mathbb R^5,qquad
\tilde{\bm q}_j=\frac{\bm q_j^{\rm meas}}{\|\bm q_j^{\rm meas}\|}.
\end{equation}
动力学为
\begin{equation}
\dot{\bm v}=\frac{1}{\mT}
\left(\frac{R(\tilde{\bm q})\ev}{\|\tilde{\bm q}\|^2}\Tphys-k_d\bm v\right)-g\ev+\bm w,qquad
\dot{\mT}=0,
\end{equation}
残差只有 $\bm e_j^v=\bm y_j^v-\bm v_j$。模型内部再次除以四元数模长平方，使非单位四元数不会重新缩放推力；由于旋转矩阵对 $q$ 为二次形式，$R(q)=R(-q)$，无需半球对齐。姿态是输入而不是状态，代价中没有四元数分量残差或模长惩罚。

位置运动学、姿态运动学和 Euler 方程本身不含质量；给定推力方向和大小后，速度响应已经包含窗口中关于质量的全部信息。全状态估计器增加的姿态轨迹反而形成转动模型误差传入质量的间接路径。代价是实测姿态噪声不再由 MHE 平滑；仿真里程计接近理想，因此这一点尚未在真实噪声下验证。

\subsection{实验批次使用的全状态估计器}

全状态估计器使用 $x_{\rm aug}=[x;\mT]\in\mathbb R^{14}$，质量在窗口内保持常值，过程噪声只作用于13个物理状态。其参数包括重构力/力矩、滞后惯量增量和固定为零的水平质心偏移。残差是位置、速度、四元数和角速度的13维分量差；质量主要仍通过速度残差辨识，转动残差则通过估计姿态间接影响质量。

\subsection{控制器预测模型}

控制器参数为
\begin{equation}
p_{\rm NMPC}=[x_{\rm ref}(13);\mT;\dJ;\cvec(2);\hat{\bm d}(3)]\in\mathbb R^{20},
\end{equation}
并使用完整平动和转动模型。水平偏移固定为零，$\hat{\bm d}$ 是残余比力。控制器是唯一需要惯量的模型，惯量由发布的质量计算而非直接估计。发布质量经过正部、低通和速率限制，不会低于空机质量。

\section{移动时域估计器}

两种估计器都求解同一类约束最小二乘问题：
\begin{align}
\min_{x_{0:N},w_{0:N-1}}\quad &
\|x_0-\bar x_0\|_{W_0}^2+
\sum_{j=0}^{N-1}\beta_j\left(\|y_j-Cx_j\|_{W_y}^2+\|w_j\|_{W_w}^2\right)\\
\text{s.t.}\quad &x_{j+1}=F_M(x_j,p_j)+Gw_j,qquad
\underline m\le\mT\le\overline m.
\end{align}
时域为 $N=20$、$\Delta t=\SI{0.1}{\second}$。每个 stage 将区间起点测量与该区间平均输入配对；最新测量是终端节点，不单独加代价。质量在窗口内没有过程噪声，从而一个窗口约束同一个质量值。

全状态估计器采用 $\sigma_v=\SI{0.05}{\meter\per\second}$ 等代表性噪声构造 Bryson 权重；过程噪声权重在平动方向取 $10^4$，迫使持续的推力--运动不一致由质量解释。质量到达代价对应 $\sigma_{m,0}=\SI{3.16}{\kg}$，很软但非零，以保持质量方向的数值曲率。降阶估计器精确复用速度、平动过程噪声和质量到达代价：
\begin{equation}
W_y^R=\sigma_v^{-2}I_3,qquad W_w^R=10^4I_3,qquad
W_0^R=\operatorname{diag}(10^3I_3,\sigma_{m,0}^{-2}).
\end{equation}

\subsection{质量边界与健康检查}

全状态估计器使用硬约束 $[0.95\mB,\SI{5.0}{\kg}]$。当推力映射低估 \SI{10}{\percent} 时，空机等效质量约为 $0.9\mB=\SI{1.858}{\kg}$，与 \SI{1.961}{\kg} 下界冲突，误差会被边界掩盖。

降阶估计器将数值盒与物理判断分开。数值盒为 $[0.70\mB,\SI{5.0}{\kg}]$，只用于远离 $1/m$ 奇点和抑制发散；物理健康区间为
\begin{equation}
\mB-\SI{0.10}{\kg}\le \hat m\le
\mB+\SI{0.30}{\kg}+\SI{0.10}{\kg},
\end{equation}
即 $[\SI{1.964}{\kg},\SI{2.464}{\kg}]$。名义推力下264架回放的降阶估计范围为 \SIrange{1.985}{2.368}{\kg}，从未越界。越界估计只被检测、不被校正；原子帧标为不健康，控制器保持上一可信模型。因为模型只含 $T/m$，平动数据无法同时估计推力比例和质量。

\subsection{输入门控、初始化与失败处理}

高度低于 \SI{0.5}{\meter} 或推力低于 \SI{1}{\newton} 时，不把重构力作为质量证据，因为地面支持力破坏 $T\approx mg$。任一输入超过 \SI{0.30}{\second} 未更新，窗口被丢弃。抓取悬停高度为 \SI{0.35}{\meter}，因此抓取阶段窗口会被清空，并在起飞后约六秒重建；抓取本身不会跨在同一个窗口中。

初值由窗口竖直力平衡获得：
\begin{equation}
m^{(0)}=
\frac{\displaystyle\sum_{i=0}^{N-1}T_i\cos\phi_i\cos\theta_i}
{\displaystyle Ng+(v_{z,N}-v_{z,0})/\Delta t}.
\end{equation}
该式只使用窗口两端速度，无需数值微分，也不包含估计器自身状态。非零求解状态或非有限解被丢弃并将帧标记为不健康；连续五次失败后重建窗口。降阶估计器只把新种子用于迭代和到达代价，发布值保持上一次成功结果，直到新的求解成功。

事件降权、指数遗忘和质量随机游走均实现过。312架释放实验中，推力阶跃触发器在81架触发，但只有10架真正改变了已求解窗口的权重；多数触发落在抓取门控间隙内。因此实际系统基本依靠均匀权重容忍一个窗口量级的参数突变，而不是精确检测变化时刻。

\section{自适应 NMPC}

NMPC 在 $N=20$、$\Delta t=\SI{0.05}{\second}$ 的预测时域上最小化跟踪误差和输入相对实时悬停基线的偏差：
\begin{equation}
T_h=\mT(g-\hat d_z),\qquad u_{\rm hover}=[T_h,0,0,0]^\top.
\end{equation}
若基线仍按空机质量计算，输入代价会持续把推力拉到真实平衡值以下，形成稳态高度误差；增大输入权重反而使误差更严重。因此基线必须随估计质量和残余比力实时更新。

姿态跟踪误差使用相对四元数 $q_{\rm ref}^*\otimes q$ 的向量部分，而不是四个分量直接相减。各残差按代表性允许误差无量纲化。控制器以 acados、显式 RK4、多重射击、Gauss--Newton、部分凝聚 HPIPM 和带回溯的完整 SQP 求解；实时迭代方案在该配置中曾静默发散，故未采用。

\subsection{残余比力}

质量估计在机动中偏低时，预测模型的推力前馈也偏低。控制器用 $\mathcal L_1$ 型状态预测器估计加性比力 $\hat{\bm d}$，经低通后加入平动模型和悬停基线。它只补偿质量误差对跟踪的后果，不进入 MHE、无载荷证据或生命周期决策。正常释放接受时将其清零；无指令脱落时由自身动态逐步消退。

\subsection{输入约束}

归一化油门 $\eta$ 与转速满足 $\omega=\omega_{\min}+\Delta\omega\eta$，推力为 $T=4k_f\omega^2$。对 $\eta\in[0.05,0.95]$，实际总推力包线为
\begin{equation}
T\in[\SI{1.27}{\newton},\SI{31.35}{\newton}].
\end{equation}
旧的按机体重量设置的 \SI{40.5}{\newton} 上界高估执行器能力 \SI{29}{\percent}。滚转/俯仰力矩上界 \SI{0.5}{\newton\meter} 较保守，偏航上界 \SI{0.2}{\newton\meter} 接近甚至略高于实际能力，这与实测偏航力矩不足一致。

\section{连续载荷接口与内环调度}

估计器每次发布一个原子帧，包含复合质量、由其计算的三轴惯量信息、无载荷证据 $\gamma$、健康位和最近成功解年龄。控制器还检查帧接收年龄。任何不健康或过期帧都冻结上一可信模型，并重置释放确认的持续计时；估计退化最多延迟确认，不能自行触发确认。

有效参数先经过一阶低通和独立速率限制。控制器侧时间常数为 \SI{0.20}{\second}，质量变化率限制为 \SI{0.60}{\kg\per\second}，惯量增量限制为 \SI{0.080}{\kg\meter\squared\per\second}。低通抑制噪声，速率限制约束相邻 NMPC 求解之间模型变化的速度，保护热启动。

载荷使 PX4 空机速率环的有效带宽按惯量比下降。系统不在线改写 PX4 增益，而是围绕当前实测角速度放大滚转/俯仰角速度误差：
\begin{equation}
\bm\omega'_{\rm cmd}=\bm\omega+\varsigma(\bm\omega_{\rm cmd}-\bm\omega),
\end{equation}
其中请求比例为惯量比 $\kappa$，上限5.0，并经过死区、\SI{0.5}{\second} 低通和非对称速率限制。基于速率余量的限制保证缩放后指令仍位于硬上限的 \SI{95}{\percent} 内，且 $\varsigma\ge1$，调度器只能增加而不能削弱原始控制权威。

抓取初期在 MHE 尚未形成载荷证据时最危险。系统在抓取使能后，将惯量参数暂时置于 \SI{0.30}{\kg} 载荷包线对应的下限，即 $\dJ=\SI{0.0191}{\kg\meter\squared}$、$\kappa=2.34$；质量通道不被人为抬高。评估配置把该惯量保护保持到释放命令，再用 \SI{0.5}{\second} 淡出到估计值。该命令只结束保护叠加层，并不宣称载荷已经离开。已知缺点是漏抓时空机会带着重载惯量保护飞行。

\section{抓取--释放生命周期}

连续参数适应与释放接受是两件事。质量、惯量和内环比例始终逐帧跟随健康估计；释放接受只改变任务状态、在有释放命令时清零残余比力并宣布任务完成，不直接设置质量或惯量。

质量证据定义为滤波载荷质量的平滑空载分数：当估计载荷不超过 \SI{0.015}{\kg} 时得分为1，不低于 \SI{0.060}{\kg} 时为0，中间采用三次 smooth-step，再经过 \SI{0.25}{\second} 低通。它是确定性证据分数，不是后验概率。释放接受条件为
\begin{equation}
\gamma\ge0.90\quad\text{在健康、新鲜帧上连续保持}\quad\SI{1.0}{\second}.
\end{equation}
若释放命令后 \SI{12}{\second} 仍未确认，系统进入 \textsc{unresolved}。该状态不是接受释放：不清模型、不清残余力，而是用 \SI{3}{\second} 五次参考轨迹制动到悬停，降低风险并给估计器提供安静工况。

\subsection{力平衡证据与双重启用}

为了避免质量 MHE 在机动中的偏差，系统另外从实测总推力和机体 $z$ 轴比力计算
\begin{equation}
m_{\rm acc}=\frac{\bar T}{\max(\bar f_z,\SI{1}{\meter\per\second\squared})}.
\end{equation}
推力和比力均用 \SI{0.5}{\second} 一阶低通，随后用与质量证据相同的阈值生成 $\gamma_{\rm acc}$。这一证据不经过 MHE 的姿态轨迹、速度历史或转动模型，即使 MHE 帧不健康仍可使用；但它同样一比一继承推力系数误差。

证据始终计算，确认只在两种情况下计数：存在待处理的正常释放命令；或意外脱落看门狗已经武装。看门狗在离地后，检测到 $\gamma_{\rm acc}<0.75$ 的有载证据连续 \SI{2}{\second} 时武装，并保持到下一次抓取。地面支撑阶段两种启用条件都不成立，因此载荷已吸附但仍由地面支撑时，不会把“空机力平衡”误认为释放。

无释放命令但看门狗已武装时的确认被声明为意外脱落：控制器清零载荷几何、结束惯量保护、恢复空机内环比例并取消之后的计划释放。系统不发送夹爪命令，因为无法判断是载荷掉落还是夹爪故障。

\begin{figure}[H]
\centering
\includegraphics[width=0.88\textwidth]{Figures/fig_lifecycle_fsm_massonly.pdf}
\caption{载荷存在状态机与控制器释放生命周期。图中文字沿用英文源图。}
\end{figure}

\section{参考轨迹与在线执行顺序}

实验使用三维 Gerono 八字轨迹
\begin{equation}
x_r=r\sin a,\qquad y_r=\tfrac r2\sin2a,\qquad
z_r=z_h+d_z\sin a,\qquad a=wt_c,
\end{equation}
偏航沿速度方向。轨迹交叉点不是尖点，能够在紧凑空间内同时激励平移和偏航。位置幅值和航向用五次 smooth-step 渐入，避免初始偏航力矩饱和。

每个控制周期依次完成：接收并转换里程计；从旋翼转速重构物理力；维护 MHE 窗口并在必要时重新播种；求解质量并执行健康检查；更新存在状态和空载证据；发布原子帧；控制器校验、滤波并更新模型和悬停基线；判断确认或 \textsc{unresolved}；热启动求解 NMPC；以 \SI{50}{\hertz} 发布推力和经内环调度的角速度指令；记录估计、残差、约束、求解器与生命周期事件。

\section{本章小结}

部署系统把载荷建模压缩为一个复合质量决策变量：降阶 MHE 只积分平动速度，实测姿态作为已知输入；NMPC 则使用完整飞行器模型和由质量派生的惯量。健康门控的原子帧同时驱动预测模型、悬停基线和内环调度。释放接受与连续参数适应分离，并可由质量证据或独立力平衡证据驱动。结构上已经明确的局限是：质量证据继承 MHE 在机动中的失真，力平衡证据则继承推力标定误差。

% =====================================================================
\chapter{实现与实验评估}

本章在软件在环环境中评估完整接口并刻画失效边界。证据包括：120架冻结随机区组 cohort、两个较小消融批次、312架力平衡释放实验以及只用于解释机制的开发实验。不同批次不合并统计。早期偏心载荷配置的数据不作为当前主要结果。

\section{平台、任务与实时性能}

系统由 Gazebo Harmonic、PX4 SITL、ROS~2 MHE/NMPC 节点、可分离关节夹爪、生命周期日志和 $\mathcal L_1$ 残余比力模块组成。所有载荷变化都由真实第二刚体的连接或分离产生，同时改变质量和惯量。20个批次中，重构滚转/俯仰力矩与计划力矩的中位相关系数为0.93/0.70、斜率为1.08/1.10；实际偏航力矩只有计划值的0.32，再次说明估计器应使用实测旋翼力而不是计划力矩。

平台空机质量为 \SI{2.0643}{\kg}，空机惯量为 $\operatorname{diag}(0.0142,0.0142,0.0210)\,\si{\kg\meter\squared}$。载荷为边长 \SI{0.15}{\meter}、质量 \SI{0.15}{\kg} 或 \SI{0.30}{\kg} 的立方体，抓取水平容差 \SI{0.03}{\meter}。任务依次为接近、NMPC 接管、对准、下降、吸附、\SI{8.4}{\second} 抬升到 \SI{6}{\meter}、半径 \SI{10}{\meter} 的八字搬运、释放和恢复。速度约为2或 \SI{4}{\meter\per\second}，释放统一发生在八字端点。

同一桌面 CPU 核上，13状态 NMPC 的求解时间中位数为 \SI{1.3}{\milli\second}，99百分位为 \SI{2.2}{\milli\second}。全状态 MHE 的2秒窗口平均求解时间中位数为 \SI{2.8}{\milli\second}、90百分位 \SI{5.4}{\milli\second}、99百分位 \SI{14.3}{\milli\second}，记录到的最大单次求解为 \SI{77.3}{\milli\second}，未超过 \SI{100}{\milli\second} 周期。本文不据此声称嵌入式硬件上也有相同余量。

\section{评估配置与实验问题}

120架 cohort 和312架释放实验都使用14状态全状态 MHE、\SI{10}{\hertz}、$N=20$、滞后惯量、窗口力平衡初值、质量证据，以及13状态 NMPC。额外固定启用残余比力、保持到释放命令的惯量保护、相对载荷对准、平滑初始航向和缩短后的 \SI{0.27}{\meter} 力臂。所有抓取居中，水平偏移补偿关闭。后续入口已切换到降阶估计器，但前述批次结果描述的是实验时的全状态系统。

需要注意：当前工作稿的 cohort 数值来自较早的“质量分数乘以由同一质量派生的惯量分数”版本。离线重评分显示71个确认中68个时刻不变、3个提前 \SI{0.1}{\second}，但最终提交前仍应以质量单分数闭环重跑替换这些历史表格。

五个实验问题分别是：质量是否能在抓取、搬运和释放过程中辨识；物理释放能否全部确认以及延迟多少；命令先于物理分离时能否避免提前确认；始终开启和命令后开启的确认策略有何取舍；释放策略和残余比力如何影响跟踪。

cohort 比较三种释放策略：A 为命令即接受；B 为全程运行质量证据确认；C 使用相同证据，但只在释放命令后开始持续计数。三种条件为 c1（\SI{0.15}{\kg}, \SI{2}{\meter\per\second}）、c2（\SI{0.15}{\kg}, \SI{4}{\meter\per\second}）和 c3（\SI{0.30}{\kg}, \SI{4}{\meter\per\second}）；场景为正常释放 d0 与命令后延迟4秒物理分离 d4。

120架按固定随机种子分8个区组运行，每格8架，不补跑。118架满足“真实吸附且存活到释放命令”的预注册资格；1架漏抓，1架仿真器在起飞前停止。所有产生飞行数据的架次均未坠毁或发散。

\section{主要 cohort 结果}

\subsection{质量估计}

载荷悬停阶段质量估计准确，逐架均值误差位于复合质量的 $-0.44$ 至 $-0.10$\%。进入八字后出现速度相关的周期性低偏：
\begin{center}\small
\begin{tabular}{lccccc}
\toprule
条件 & 悬停误差 & 八字误差 & 释放后误差 & 载荷塌陷比例 & 质量下界比例\\
\midrule
c1 & $-0.21\%$ & $-0.97\%$ & $-0.94\%$ & 0.00 & 0.000\\
c2 & $-0.22\%$ & $-4.09\%$ & $-1.70\%$ & 0.38 & 0.010\\
c3 & $-0.29\%$ & $-3.30\%$ & $-1.68\%$ & 0.03 & 0.014\\
\bottomrule
\end{tabular}
\end{center}
\SI{4}{\meter\per\second} 时误差约每11秒下降一次，即八字周期的一半，首次加速进入八字时中位轨迹会触及质量下界。相同复合质量相对误差对轻载意味着更大的载荷相对误差，因此 c2 的“估计载荷低于真实载荷一半”比例最高。惯量保护和残余比力避免了控制失效，但没有消除质量塌陷及其对生命周期判断的影响。

\subsection{释放完成、延迟与提前确认}

118架合格飞行全部完成释放，没有进入 \textsc{unresolved}。B策略48/48完成，确认发生在物理分离后 \SIrange{2.3}{3.2}{\second}；约1秒来自持续计时，其余来自质量估计和输出低通下降到阈值。

在4秒延迟分离场景中，A策略24/24都在载荷仍连接时接受，约提前4秒；B、C共47架没有一次提前接受，物理分离后的确认延迟与正常场景相同。这证明的是生命周期语义正确性，不是跟踪性能改善，因为控制模型仍连续跟随质量估计。

B策略在高速载荷搬运阶段发生严重误确认：c2中14/16架、共78次；c3中6/16架、共36次；c1没有发生。受影响架次通常约每11秒确认一次，与质量周期塌陷同步。C策略23架中没有一次搬运期确认，也没有增加真实分离后的确认延迟。因此，对质量证据而言，证据应在命令后才允许确认；代价是无法检测无命令脱落。

正常释放中，A与B的释放后位置 RMSE 没有一致差异。搬运段 RMSE 主要由速度决定：\SI{2}{\meter\per\second} 约 \SI{0.03}{\meter}，\SI{4}{\meter\per\second} 约 \SI{0.12}{\meter}。确认策略不改变连续质量模型，所以不应把B写成跟踪更好。

\section{312架力平衡释放实验}

第二个冻结批次于2026年11月3--4日运行，所有312架均合格且无发散。它使用双重启用的力平衡规则 D，并增加 c4（\SI{0.30}{\kg}, \SI{2}{\meter\per\second}）。四组实验为：E1 A/D与正常/延迟分离，144架；E2 Fixed/$\mathcal L_1$/MHE控制器，72架；E3 检测器推力映射0.9/1.1，48架；E4无命令脱落/夹爪卡死，48架。

\subsection{释放语义}

A策略在延迟分离时36/36提前确认，中位提前 \SI{4.05}{\second}。D策略72/72完成，延迟分离时0/36提前确认；确认位于物理分离后 \SIrange{1.90}{2.51}{\second}，中位 \SI{2.20}{\second}、95百分位 \SI{2.44}{\second}。重载使延迟增加约 \SI{0.30}{--}\SI{0.35}{\second}，高速影响仅约 $-0.08$ 至 $-0.14$ 秒，质量--速度交互不显著。

在240架采用D生命周期的飞行中，看门狗在228架武装；从武装到释放命令的整个搬运阶段没有一次误确认。未武装的12架全部来自推力低估轻载工况。D与A跟踪总体无一致差异，仅重载高速、延迟分离一格中D的释放后 RMSE 中位高 \SI{0.014}{\meter}。

\subsection{无命令脱落与夹爪卡死}

24架无命令脱落中，看门狗全部在脱落前武装，24/24检测成功，延迟 \SIrange{1.8}{2.5}{\second}，没有提前声明。每架都清零载荷几何、恢复空机内环比例、取消之后的计划释放且没有发散。

24架夹爪卡死中，0/24误接受，载荷模型从未被清除。全部在12秒后进入 \textsc{unresolved}，执行制动并安全悬停，所有架次悬停阶段最大速度低于 \SI{0.18}{\meter\per\second}。

\subsection{推力映射误差}

力平衡质量一比一继承推力比例误差。推力高估 \SI{10}{\percent} 时，释放后的空机仍显得过重，24/24都漏检并进入 \textsc{unresolved}，随后安全悬停。推力低估 \SI{10}{\percent} 时，48架中相应24架全部完成，但轻载的12架在搬运阶段得分始终为1.0，导致意外脱落看门狗12/12无法武装；若再叠加延迟分离，预计可能在真实分离前接受，但该组合没有实测。重载时看门狗能够武装，但一架搬运期得分达到0.85，接近0.90阈值。

\subsection{Fixed、\texorpdfstring{$\mathcal L_1$}{L1} 与 MHE控制器}

三种控制器的确认延迟无明显差异。相对固定参数控制器，MHE在24/24配对中搬运 RMSE 更小，中位改善 \SIrange{0.04}{0.14}{\meter}。相对使用相同残余比力但固定质量的 $\mathcal L_1$ 控制器，搬运跟踪没有一致差异；c2反而有5/6配对由 $\mathcal L_1$ 更好。释放后，MHE只在两个低速条件显示 \SIrange{0.006}{0.017}{\meter} 的较小峰值优势。因三者共享惯量保护和内环调度，E2不检验载荷模型的转动收益。

\section{消融实验}

\subsection{残余比力}

48架冻结消融比较残余比力开/关，每个条件8个配对。全部吸附、释放完成且无发散。开启残余比力后，搬运 RMSE 在 c1/c2/c3 分别从0.047/0.148/0.145米降至0.033/0.116/0.106米，下降 \SIrange{22}{31}{\percent}；平均竖直误差从 $-4.5/-28.4/-44.3$ 毫米变为 $+0.4/-0.8/+0.2$ 毫米；释放后 RMSE 下降 \SIrange{33}{48}{\percent}。它没有改变质量偏差、确认延迟或搬运期误确认，只补偿质量偏差对跟踪的后果。

\subsection{估计器内部的惯量处理}

让质量实时改变估计器转动惯量的“实时”方案在单架开发实验中造成28--64\%质量高估甚至发散，因此弃用。16架滞后/省略惯量对比中，搬运位置误差0.121米对0.113米，无显著改善；但质量偏差由 $-8.9\%$ 减至 $-4.0\%$，16架交叉回放全部由滞后方案更准确。因此全状态配置采用滞后惯量：它约将机动质量偏差减半，但不消除偏差，也没有显示跟踪收益。

\subsection{全状态与降阶估计器}

对312架释放实验日志进行降阶回放，在264架名义推力飞行中得到：
\begin{center}\small
\begin{tabular}{lrrrr}
\toprule
指标 & F在线 & R回放 & R更优 & $p$\\
\midrule
带载质量 MAE [kg] & 0.0359 & 0.0127 & 264/264 & $5\times10^{-45}$\\
释放后质量 MAE [kg] & 0.0291 & 0.0093 & 240/240 & $4\times10^{-41}$\\
离地后收敛 [s] & 4.21 & 4.22 & -- & --\\
架次内回到空机值 & 183/240 & 240/240 & -- & --\\
中位求解时间 [ms] & 2.8 & 0.16 & 264/264 & --\\
\bottomrule
\end{tabular}
\end{center}
R在每一个带载和释放后窗口都更准确；57条F轨迹在高速持续八字中稳定于空机值下方约2.15\%，因而在架次结束前没有重新进入 $\pm\SI{0.03}{\kg}$ 区间。回放不能评价闭环跟踪，使用近理想仿真里程计，也没有真实复现抓取期间约6秒的窗口间隙。

使用F的硬下界回放推力误差时，$-10\%$ 工况22/24架出现共107次最大迭代，均在质量下界附近；这促使部署的R改用宽数值盒和独立健康检查。四架闭环冒烟中没有非零求解状态或重初始化。$-10\%$ 两架在空机时被健康检查识别，控制器保持模型且正常确认；$+10\%$ 两架均进入 \textsc{unresolved} 并安全悬停。重载 $+10\%$ 时估计约 \SI{2.60}{\kg}，超出健康带，控制器却只能保持越界瞬间的最后可信值 \SI{2.272}{\kg}；健康检查避免不合理模型，但不能恢复正确质量。

当前英文工作稿已经预注册F/R闭环配对：2种质量、2种速度、立即/延迟4秒分离、每格6对，共96架；D释放规则、控制器、轨迹、几何和安全设置固定。\textbf{截至本中文稿生成时，英文源文件仍标注该批次结果待填，不能把回放结果写成闭环结果。}

\subsection{水平偏移与学习权重}

早期偏心抓取消融在 \SI{0.15}{\kg}、名义 \SI{0.10}{\meter} 偏移下比较4对飞行。建模偏移与不建模的搬运 RMS 为0.1230米和0.1213米，未显示收益；日志表明约99\%配平力矩由 PX4 速率环积分项承担。这不是大偏心等效性证明，当前系统通过抓取前对准避开该区域。

参数化阶段权重经交叉熵搜索后，在配对消融中没有优于手工规则；目标在手工规则附近较平坦。真正有用的是识别模型已经变化并移除陈旧数据，而不是精确学习触发后的权重幅值。

\section{失效模式与工作边界}

主要失效是全状态质量估计在高速八字中周期性向空机值塌陷。早期开发批次显示：相对对准和航向平滑提高生存率但未消除塌陷；缩短力臂把塌陷比例从0.05--0.90降到0.03--0.24；16架力臂对比中，塌陷比例与实测臂长的 Spearman 相关为0.66。滞后惯量减小偏差，惯量保护和残余比力分别保护内环与平动跟踪，命令后武装避免搬运误确认，但这些措施都没有修复全状态估计的根因。

重载高速是压力条件：悬停需要最大推力的 \SI{74}{\percent}，名义惯量比2.34，实际臂长范围下最大约2.7，仍低于角速度缩放上限5.0。早期 \SI{0.47}{\meter} 力臂时惯量比5.07，超过调度能力，因此缩短夹爪。当前 cohort 中滚转力矩从未到界，俯仰平均最多0.3\%样本到界，高速偏航约2.5\%样本到界。

已知结构缺口是相位定时释放命令没有由估计健康度门控。若估计被冻结时仍执行释放，物理分离后控制器可能继续持有旧模型。另一个配置风险是入口默认与节点库默认不同：实验入口选择质量单参数形式并已切换到降阶估计器，而直接启动节点仍可能进入论文未使用的16状态一阶矩形式。每个批次必须记录全部解析后开关。

\section{有效性威胁、缺失基线与复现}

全部系统级证据来自软件在环、单一机体和夹爪、居中刚性抓取、无风、近理想里程计。实际力臂逐架变化但没有逐架测量；抓取柔性、摆动、速度噪声、传输延迟和真实诱导流均未覆盖。持续竖直风与载荷变化在只看总推力的残差上本质不可区分。

推力标定误差会一比一传入质量：对约 \SI{2.2}{\kg} 复合质量，$\pm5\%$ 即 $\pm\SI{0.11}{\kg}$，远大于 \SI{0.024}{\kg} 确认阈值。低估方向可能把真实载荷看成空载，高估方向可能在释放后留下虚假载荷；后者进入 \textsc{unresolved} 较安全，前者可能静默失去保护。检测器的 $\pm10\%$ 实验已经验证两种方向性失败，但估计器侧的 $\pm5\%$ 注入、部分载荷逐步丢失和窗口初值对比仍未正式运行。

已有固定参数和 $\mathcal L_1$ 基线，但没有 INDI、完整旋转通道的 $\mathcal L_1$-NMPC，或在相同误报约束下优化的 CUSUM 变化检测器。对释放问题而言，缺少优化检测器是最大的剩余基线缺口。

每次运行记录代码版本、未提交补丁、未跟踪文件归档、全部环境开关、问题维数、时域、权重、边界、事件时间、随机种子和求解器统计。数值失败保留在分母中。批次并非来自干净提交，因此论文只能声称可以由“版本+补丁+归档”精确重建，不能声称由某一个提交单独产生。

\section{本章小结}

120架 cohort 表明：全状态估计器在带载悬停准确，但在 \SI{4}{\meter\per\second} 持续机动中周期性低偏；118架合格释放全部完成，质量证据确认位于物理分离后 \SIrange{2.3}{3.2}{\second}；证据门控避免了延迟分离时的提前接受，但全程开启的质量规则在高速搬运中误确认114次，命令后武装则没有；确认策略不改善跟踪，残余比力显著改善跟踪但不改善质量估计。

312架力平衡实验中，D规则没有搬运期误确认或延迟分离提前确认，检测24/24无命令脱落，对24/24夹爪卡死均不误接受；但推力系数误差会在两个方向产生不同失效。相对固定模型，适应在搬运中明显更好；相对仅用残余比力的 $\mathcal L_1$，显式估质量没有一致跟踪收益。

降阶估计器的312架回放在每一个名义飞行中都具有更小质量误差，四架闭环冒烟验证了宽数值盒和健康门控的基本安全行为；但96架正式闭环配对在当前英文稿中仍待填。因此当前证据支持把降阶结构作为更合理的质量估计器，同时必须把离线回放、冒烟结果和正式闭环验证严格区分。

\end{document}
